\begin{abstract}
Task-order stress tests are increasingly used to evaluate self-improving agents. The same tasks are permuted while the agent accumulates experience, so the resulting performance change is treated as sensitivity to learning history. We show that stateful web benchmarks turn this experiment into a joint intervention. In WebArena, earlier tasks mutate shared services and can change the ground truth of later tasks. A permutation therefore changes agent memory history and benchmark environment history at the same time. We audit the released trajectories of a recent study that reports task-order fragility for Agent Workflow Memory (AWM) and ReasoningBank. A direct released trajectory shows a purchase reaching checkout despite benchmark score zero. A later task reads the resulting current order correctly and then fails against stale fixed ground truth. Three preregistered dynamic-reference sentinels reproduce this mechanism across both memory methods. All 18 ordinal evaluations pass, while all 36 shuffled evaluations fail. A conservative retrospective over 84 run-domain cells identifies 887 state-exposed reader units in the shuffled memory runs; 424 fail. Removing state-sensitive readers reduces AWM's measured ordinal-to-shuffle gap by 34.6--44.5\% and ReasoningBank's by 14.9--17.5\%. The remaining gap captures residual order sensitivity after the environment-sensitive stratum is removed. These results establish environment carryover as a material confound in published task-order evaluation. We introduce a state-isolated order-stress protocol that restores service state per task and adds a matched shuffled no-memory control. Live-world evaluations can preserve evolving services through current-state effect scoring. The protocol isolates memory-history sensitivity while retaining interactive benchmark realism.
\end{abstract}
